% ============================================================
% INTERNAL CROSS-REFERENCES
% Viability Philosophy System
% Book I: Ontology of Viability
% ============================================================

\cleardoublepage

\section*{Internal Cross-References}

\phantomsection
\addcontentsline{toc}{section}{Internal Cross-References}

% ------------------------------------------------------------
% 1. PURPOSE AND SCOPE
% ------------------------------------------------------------

\subsection*{1. Purpose and Scope}

Internal cross-references establish explicit connections
between different locations and logical entities within
the Viability Philosophy System.

Their purpose is to enable the reader to move between
related passages, identify where a subject has been
introduced or developed, and trace the relations between
different components of the book.

An internal cross-reference may point to a structural
location, a logical entity, or both, depending on the
nature of the reference.

The referencing system must preserve the distinction
between the location of a textual passage and the identity
of the entity discussed in that passage.

Internal cross-references serve the following functions:

\begin{itemize}
    \item locating passages that are relevant to the
    discussion at hand;
    \item directing the reader to earlier or later
    treatments of a subject;
    \item connecting different discussions of the
    same logical entity;
    \item identifying the textual locations at which
    definitions, propositions, arguments, and other
    logical entities are introduced or examined;
    \item making relations between different parts
    of the book traceable;
    \item preserving the accuracy of references
    during manuscript revision.
\end{itemize}

This document establishes the principles governing
the construction, use, interpretation, and maintenance
of internal cross-references.

% ------------------------------------------------------------
% 2. TWO FORMS OF INTERNAL REFERENCE
% ------------------------------------------------------------

\subsection*{2. Two Forms of Internal Reference}

The internal referencing system distinguishes between
two principal forms of reference:

\begin{enumerate}
    \item structural references;
    \item logical entity references.
\end{enumerate}

These forms of reference are related but perform
different identification functions.

\subsubsection*{2.1. Structural References}

A structural reference identifies a textual location
within the hierarchical organization of a book.

It uses the structural identifier assigned to the
relevant part, chapter, section, subsection, or paragraph.

For example:

\begin{center}
    \texttt{VS.B01.P01.C01.S01.SS01.P001}
\end{center}

This identifier designates a particular paragraph
within the structure of Book I.

A structural reference is appropriate when the reader
needs to locate a passage, inspect a particular
discussion, or examine the context in which a claim
is presented.

\subsubsection*{2.2. Logical Entity References}

A logical entity reference identifies an entity
within the logical organization of the Viability
Philosophy System.

For example:

\begin{center}
    \texttt{AX-001}
\end{center}

This identifier designates a logical entity rather
than a particular textual location.

A logical entity may be discussed in several chapters,
sections, or paragraphs without acquiring a new
logical identifier at each occurrence.

A logical entity reference is therefore appropriate
when the intended object of reference is the entity
itself rather than the passage in which it appears.

% ------------------------------------------------------------
% 3. PRINCIPLE OF IDENTIFICATION
% ------------------------------------------------------------

\subsection*{3. Principle of Identification}

Every internal cross-reference must make clear
what kind of object it identifies.

A reference to a textual location must not be
presented as though it were an identifier of
a logical entity.

Likewise, a reference to a logical entity must
not be treated as though it identified a single
passage in which that entity is discussed.

The distinction can be summarized as follows:

\begin{center}
\begin{tabular}{ll}
    \toprule
    \textbf{Reference Type} & \textbf{Identifies} \\
    \midrule
    Structural reference
        & A location in the textual hierarchy \\
    Logical entity reference
        & An entity in the logical system \\
    \bottomrule
\end{tabular}
\end{center}

Where the discussion requires both forms of
identification, both should be supplied.

For example, a passage may refer to a logical
entity by its identifier and also provide the
structural location at which the entity was
introduced or formally examined.

The two references must remain separately
recognizable.

% ------------------------------------------------------------
% 4. REFERENCES TO STRUCTURAL LOCATIONS
% ------------------------------------------------------------

\subsection*{4. References to Structural Locations}

Structural cross-references should identify
the narrowest textual location necessary
for the purpose of the reference.

A reference to an entire part or chapter may
be sufficient when the discussion concerns
the general argument developed there.

A reference to a section, subsection, or
paragraph is preferable when the intended
passage is more specific.

The level of precision should correspond
to the function of the reference.

\subsubsection*{4.1. References to Parts and Chapters}

References to parts and chapters are appropriate
when directing the reader to a substantial
thematic or argumentative unit.

Such references should identify the relevant
structural unit unambiguously.

\subsubsection*{4.2. References to Sections and Subsections}

References to sections and subsections are
appropriate when the relevant discussion
is confined to a more specific portion
of a chapter.

They should be used when a chapter-level
reference would leave the intended location
insufficiently precise.

\subsubsection*{4.3. References to Paragraphs}

Paragraph-level references should be used
when the intended claim, distinction,
qualification, or explanation occurs
within a specific paragraph.

A paragraph reference must correspond
to an existing paragraph identifier.

Paragraph references must not be confused
with references to individual sentences,
typographical lines, or page positions.

% ------------------------------------------------------------
% 5. REFERENCES TO LOGICAL ENTITIES
% ------------------------------------------------------------

\subsection*{5. References to Logical Entities}

A logical entity reference should use
the identifier assigned to the entity
in the logical entity identification
system.

The identifier must be reproduced
accurately and consistently.

For example:

\begin{center}
    \texttt{AX-001}
\end{center}

A logical entity reference remains valid
across different textual occurrences
of the same entity.

The reader should not be required to
infer the identity of an entity solely
from its name or from the paragraph
in which it happens to appear.

Where an entity has an established
identifier, that identifier should
be preferred when precise identification
is necessary.

The assignment and classification
of logical entity identifiers are
governed by the document
\textit{Logical Entity Identifiers}.

% ------------------------------------------------------------
% 6. COMBINED REFERENCES
% ------------------------------------------------------------

\subsection*{6. Combined References}

Some discussions require both a reference
to a logical entity and a reference
to a structural location.

This occurs, for example, when the reader
must identify an entity and then locate
a particular passage in which that entity
is introduced, defined, interpreted,
qualified, or applied.

In such cases, the reference should
provide both forms of identification.

The logical entity identifier establishes
which entity is intended.

The structural identifier establishes
where the relevant discussion occurs.

These references must not be collapsed
into a single identifier or treated
as interchangeable.

A combined reference is particularly
useful when:

\begin{itemize}
    \item an entity is discussed in
    multiple locations;
    \item the reader must distinguish
    the entity from a particular
    interpretation of it;
    \item a later argument depends
    on a formulation presented
    in an earlier passage;
    \item the text refers to an entity
    while directing the reader
    to a specific textual treatment.
\end{itemize}

% ------------------------------------------------------------
% 7. DIRECTION AND FUNCTION OF REFERENCES
% ------------------------------------------------------------

\subsection*{7. Direction and Function of References}

Internal cross-references may direct
the reader to different locations
according to the argumentative
or explanatory function of the reference.

The principal functions include:

\subsubsection*{7.1. Forward References}

A forward reference directs the reader
to a location or entity that will be
examined later in the book.

It may be used when a discussion
introduces a question, distinction,
or problem whose full treatment
is deferred.

A forward reference should not imply
that the later discussion has already
established a conclusion that has
not yet been presented.

\subsubsection*{7.2. Backward References}

A backward reference directs the reader
to an earlier location or entity.

It may be used to recall a definition,
distinction, argument, or established
formulation on which the current
discussion depends.

A backward reference should identify
the relevant earlier material with
sufficient precision.

\subsubsection*{7.3. Lateral References}

A lateral reference directs the reader
to another location that is related
to the current discussion but does
not necessarily stand in a direct
sequence of prior establishment
and subsequent development.

Such references may connect discussions
that address different aspects of
the same problem or examine related
entities from different perspectives.

\subsubsection*{7.4. Referential Function}

The presence of a cross-reference
does not, by itself, establish the
logical nature of the relation
between the passages or entities
being connected.

A reference may indicate that
a subject is discussed elsewhere,
but it does not automatically
establish dependence, derivation,
equivalence, contradiction, or
logical entailment.

When the nature of a logical relation
matters, it must be stated or
specified through the appropriate
logical dependency system.

% ------------------------------------------------------------
% 8. ACCURACY AND REFERENTIAL PRECISION
% ------------------------------------------------------------

\subsection*{8. Accuracy and Referential Precision}

Every internal cross-reference
must satisfy the requirements
of accuracy and referential precision.

A reference is accurate when it
points to an existing structural
location or a valid logical entity
identifier.

A reference is precise when it
identifies the intended object
without unnecessary ambiguity.

The following principles apply:

\begin{itemize}
    \item A structural reference must
    correspond to the actual hierarchy
    of the manuscript.

    \item A logical entity reference
    must correspond to the entity
    assigned that identifier.

    \item A reference must not point
    to an obsolete location when
    the intended passage has been
    relocated.

    \item A reference must not use
    an identifier belonging to
    a different entity.

    \item The level of structural
    specificity must be appropriate
    to the claim or discussion
    being referenced.

    \item Where several locations
    are relevant, the reference
    must make clear which location
    supports which part of the
    discussion.
\end{itemize}

% ------------------------------------------------------------
% 9. REFERENCES AND ARGUMENTATIVE CONTEXT
% ------------------------------------------------------------

\subsection*{9. References and Argumentative Context}

An internal cross-reference should
serve a discernible function within
the surrounding argument.

It should help the reader understand
why another passage or entity is
relevant to the present discussion.

A reference should not be used as
a substitute for explaining a
necessary premise, distinction,
or argumentative transition.

When the current discussion depends
on a previous formulation, the
relevant relation should be stated
clearly enough for the reader
to understand the dependency.

Similarly, when a later discussion
develops or qualifies an earlier
claim, the cross-reference should
help the reader trace that development.

The purpose of cross-referencing
is not merely to increase the number
of connections within the text.
It is to make significant relations
between textual and logical elements
traceable.

% ------------------------------------------------------------
% 10. MAINTENANCE DURING REVISION
% ------------------------------------------------------------

\subsection*{10. Maintenance During Revision}

Internal cross-references must be
reviewed whenever the manuscript's
structure or logical inventory
is revised.

A change in one location may affect
references elsewhere in the book.

The following cases require review.

\subsubsection*{10.1. Structural Relocation}

When a referenced passage is moved
to another structural location,
the corresponding structural
references must be updated.

The logical entity identifiers
associated with that passage
must remain unchanged unless
the identity of the entities
themselves has changed.

\subsubsection*{10.2. Structural Renumbering}

When structural numbering changes,
all references that depend on
the affected identifiers must
be checked.

No reference should continue
to point to a location that
no longer corresponds to the
intended passage.

\subsubsection*{10.3. Logical Entity Revision}

When a logical entity is revised,
its identifier must be examined
according to the rules governing
the preservation of logical
entity identity.

If the entity remains the same,
its identifier must be retained.

If it is replaced by a distinct
entity, the logical entity
inventory and all relevant
cross-references must be reviewed.

\subsubsection*{10.4. Deletion of Referenced Material}

When a referenced passage or
entity is removed, references
to it must be reviewed.

They must be corrected, redirected,
or removed according to the
actual function of the reference.

A reference must not be retained
merely because it existed in
an earlier version of the manuscript.

% ------------------------------------------------------------
% 11. INTERNAL REFERENCES AND
% EXTERNAL SOURCES
% ------------------------------------------------------------

\subsection*{11. Internal References and External Sources}

Internal cross-references identify
locations and entities belonging
to the Viability Philosophy System.

They must be distinguished from
bibliographic references to external
books, articles, authors, documents,
and other sources.

An external source does not become
an internal structural element
or logical entity merely because
it is cited in the book.

Where an internal discussion refers
to an external source, the two
forms of reference may appear
together, but they retain their
separate functions.

The external source is identified
through the bibliographic system.

The internal passage is identified
through its structural identifier.

Any internal logical entity
discussed in relation to that
source retains its own logical
entity identifier, where one
has been assigned.

% ------------------------------------------------------------
% 12. RELATION TO THE OTHER
% REFERENCING DOCUMENTS
% ------------------------------------------------------------

\subsection*{12. Relation to the Other
Referencing Documents}

Internal cross-referencing forms
one component of the broader
referencing system.

Its operation must remain consistent
with the other components of
that system.

In particular:

\begin{itemize}
    \item \textit{Structural Identifiers}
    establish the identification
    of textual locations.

    \item \textit{Logical Entity
    Identifiers} establish the
    identification of logical
    entities.

    \item \textit{Entity Location
    References} address the relation
    between entities and their
    textual locations.

    \item \textit{Logical Dependencies}
    specify relevant relations
    and dependencies between
    logical entities.

    \item \textit{External Source
    References} govern references
    to sources outside the system.

    \item \textit{Bibliographic
    Organization} establishes
    the organization of bibliographic
    information.
\end{itemize}

These components must be coordinated
without collapsing their distinct
identification and referencing
functions.

% ------------------------------------------------------------
% 13. SUMMARY
% ------------------------------------------------------------

\subsection*{13. Summary}

Internal cross-references connect
textual locations and logical
entities within the Viability
Philosophy System.

They operate through two distinct
forms of identification:

\begin{enumerate}
    \item structural identifiers,
    which identify locations within
    the textual hierarchy;

    \item logical entity identifiers,
    which identify entities within
    the logical organization
    of the system.
\end{enumerate}

A reference may use either form
or both, depending on its purpose.

All internal cross-references
must be accurate, sufficiently
precise, and maintained throughout
the revision of the manuscript.

They must also distinguish
referential connection from
logical dependency: pointing
to another passage or entity
does not, by itself, establish
the logical relation between them.

The system of internal
cross-references therefore
provides a traceable connection
between the textual architecture
and the logical organization
of the book.

\cleardoublepage

% ============================================================
% END OF INTERNAL CROSS-REFERENCES
% ============================================================